\documentclass[a4paper, 12pt]{article}

% золотой набор преамблы
\usepackage[english, russian]{babel} % Подключает языковые модули. Настраивает правила переноса слов и автоматически переводит системные слова (например, вместо "Contents" пишет "Оглавление", вместо "Theorem" — "Теорема"). Последний язык в списке (russian) становится основным.

\usepackage[T2A]{fontenc} % Задает внутреннюю кодировку шрифтов LaTeX. T2A — это стандартная кодировка для поддержки кириллицы. Благодаря ей в итоговом PDF-файле русский текст будет нормально копироваться и искаться.

\usepackage[utf8]{inputenc} % Говорит компилятору, что твой исходный код .tex написан в кодировке UTF-8. Как я говорил ранее, в современном Overleaf этот пакет уже встроен по умолчанию, его можно опускать.

\usepackage{amsmath, amsfonts, amssymb, amsthm, mathtools} % интерфейс работы с теоремами обеспечивается пакетом amsthm
\usepackage{graphicx} % это главный пакет в LaTeX для вставки и редактирования изображений.
\usepackage[left=2cm, right=2cm, top=2cm, bottom=2cm]{geometry}
\usepackage{hyperref} % работа с ссылками, делает кликабельным
\hypersetup{
    colorlinks=true,    % Включает цветной текст ссылок вместо уродливых рамок
    linkcolor=blue,     % Цвет внутренних ссылок (оглавление, формулы)
    urlcolor=cyan,      % Цвет внешних веб-ссылок
    citecolor=green,    % Цвет ссылок на литературу
    pdftitle={Конспект по матанализу}, % Метаданные для PDF-файла
    pdfauthor={Студент 1-го курса}
}

\usepackage{indentfirst} % пакет для абзаца каждой секции
\usepackage{titleps} % пакент для колонтитулов
\usepackage{soulutf8} % пакет для разного вида шрифта
\usepackage{xcolor} % для работы с цветами

% для работы с таблицами
\usepackage{array}
\usepackage{tabularx}
\usepackage{multirow}

% стиль страницы
\newpagestyle{main}{ % main - название стиля, и в других фигурных скобках это содержимое
    \setheadrule{0.4pt} % толщина линии отделяющий хедер или футер от текста
    \sethead{ФПМИ Конспект}{}{Глава \thesection} % верхний колонтитул (header) обязательные три параметра!
    % \thesection — выведет номер раздела (например, 1.2)
    % \sectiontitle — автоматически выведет имя текущего раздела
    \setfootrule{0.4pt}  % толщина линии отделяющий хедер или футер от текста
    \setfoot{Конспект лекций}{\thepage}{2026} % нижний колонтитул (footer)
}
\pagestyle{main} % команда для приминения стиля страницы, можно использовать в каком-то определенном промежутке кода

% стиль plain
\theoremstyle{plain} % модификатор, который задает стиль всем следуюшим теоремам, пока не переопределить.
\newtheorem{theorem}{Теорема}
\newtheorem{lemma}{Лемма}
\newtheorem{axeom}{Аксиома}
\newtheorem{corollary}{Следствие}[theorem] % последний необязательный аргумент, говорит о нумерации следствия данной теоремы.

% стиль defenition
\theoremstyle{definition}
\newtheorem{definition}{Определение}[theorem] % привязанность к данной теореме
\newtheorem{property}{Свойство}[theorem]
\newtheorem{exercise}{Задача}[section]

% стиль remark
\theoremstyle{remark}
\newtheorem{remark}[theorem]{Замечание}

\begin{document}
    \section{\centering Мой 1 раздел}
        section выступает в роли заголовок уровня h1 как в html
        
        \subsection{Мой подраздел}
            subsection выступает в роли заголовка следующего уровня h2
        
        \subsection*{Подраздел, но без номер}
            Если стоит \* это значит что подраздел не пронумерован
            Если бы заголовок h1 был бы тоже без номера то подзаголовки брались бы от нулевого уровня
        
        \subsection{2 Подраздел}
            А вот и второй подраздел, он автоматически пронумеруется
            
            \subsubsection{Под подзаголовок}
                Этот уже уровня h3
                Секции может быть сколь-угодно много, я привел только один пример
    
    
    \section{\centering Стиль Текста}
        Тут будет показано как работать со стилем текста
        
        \subsection{Использование команд}
            \textbf{Жирный}, \textit{Курсив}, обычный
        
        \subsection{Использование областей видимости и модификаторов}
        {\bfseries \itshape Здесь выделенное}, а здесь нет. Тут у меня первая команда делает текст жирным, а вторая придает тексту форму курсива
    
    
    \section{ \centering Размер и шрифт}
        Изменение размера шрифта производится с помощью модификаторов.
        
        \subsection{Большой и маленький текст}
        {\Large Большой текст}, текст поменьше, {\small вообще маленький текст}
        
        \subsection{Доп теория}
            Заменять шрифт на один из стандартных вариантов можно как модификатором, так и командой
    
    
    \section{Окружения}
        Окружение это способ инкапсулировать много модификаторов стиля в единый блок, в который
        потом можно будет заворачивать текст. Использование окружения имеет вид
        \text{\backslash begin\{enviroment\}... \backslash end\{enviroment\} }
        
        \subsection{Среда-центр}
            \begin{center}
                Все что здесь, даже секшены будут помещены в центр, полезно для картинок, таблиц и так далее
            \end{center}
        
        \subsection{Выравнивание по правому краю}
            \begin{flushright}
                здесь будет выравнивание по центру, опять таки, касается всего что в этом в блоке
            \end{flushright}
        
        \subsection{Выравнивание по левому краю}
            \begin{flushleft}
                здесь выравнивание по левому краю, тоже касательно всего что в этом блоке, может оба эти метода когда то мне пригодятся
            \end{flushleft}
    
    
    \section{ \centering Окружения теоремы}
        Особый случай окружения это теорема. Окружение такого типа предваряет текст внутри надписью вида "Теорема N" с согласованной нумерацией.
        
        \subsection{Простейший вид теоремы}
            \begin{lemma}
                1 + 1 = 3.
            \end{lemma}
            \begin{theorem}
                2 + 2 = 5.
                $2 + 2 = 5$.
            \end{theorem}

\end{document}